\documentclass[main.tex]{subfiles}


\begin{document}

%from pagenumber to up
% \phantomsection 
% \hypertarget{toc}{}

 

 

% номер секции вручную
\setcounter{section}{33
}
\addtocounter{section}{-1} 

\section{Магнитное поле кольца с током}


Пусть имеется провод в форме кольца, через который протекает электрический ток. На рис.\,\ref{fig:p149} показана картина линий магнитного поля кольца с током.

\makeatletter

\pgfkeys{/pgf/decoration/.cd,
         start color/.store in =\startcolor,
         end color/.store in   =\endcolor
}

\pgfdeclaredecoration{width and color change}{initial}{
 \state{initial}[width=0pt, next state=line, persistent precomputation={%
   \pgfmathdivide{50}{\pgfdecoratedpathlength}%
   \let\increment=\pgfmathresult%
   \def\x{0}%
 }]{}
 \state{line}[width=.5pt,   persistent postcomputation={%
     \pgfmathadd@{\x}{\increment}%
     \let\x=\pgfmathresult%
   }]{%
   \pgfsetlinewidth{\x/150*0.035pt+\pgflinewidth}%
   \pgfsetarrows{-}%
   \pgfpathmoveto{\pgfpointorigin}%
   \pgfpathlineto{\pgfqpoint{.75pt}{0pt}}%
   \pgfsetstrokecolor{\endcolor!\x!\startcolor}%
   \pgfusepath{stroke}%
 }
 \state{final}{%
   \pgfsetlinewidth{\pgflinewidth}%
   \pgfpathmoveto{\pgfpointorigin}%
   \color{\endcolor!\x!\startcolor}%
   \pgfusepath{stroke}% 
 }
}

\makeatother


    \begin{figure}[H]
    \centering

    \begin{tikzpicture}[font=\small,            line cap=round,                             >={Stealth[bend,inset=0pt,round,length=3mm, angle'=20]},vec/.style={>={Stealth[inset=0pt,round,length=3.5mm, angle'=20]}}]


\def\lm{1} 
\def\xr{.25}


%%%%%%%%%%%%%%%%%%%%%%%%%%%%%
%%%%%%%%%%%%%%%%%%%%%%%%%%%%%

    
%lines left
%1
\def\ds{.4}
\path[yshift=.75cm,->] (0.25,0) 
arc [start angle=-90, end angle=0, x radius={\ds*1cm}, y radius={\ds*1cm-.02cm}]
arc [start angle=0, end angle=90, x radius={\ds*1cm}, y radius={\ds*1cm+.05cm}]
coordinate (e1)
;

\draw[gray,yshift=.75cm] (e1) 
arc [start angle=90, end angle=180, x radius={\ds*1cm}, y radius={\ds*1cm+.05cm}] 
arc [start angle=180, end angle=270, x radius={\ds*1cm}, y radius={\ds*1cm-.02cm}] 
;


%2
\def\df{.9}
\path[yshift=.5cm,->] (0.25,0) 
arc [start angle=-90, end angle=0, x radius={\df*1cm}, y radius={\df*1cm-.1cm}]
arc [start angle=0, end angle=90, x radius={\df*1cm}, y radius={\df*1cm+.1cm}]
coordinate (e1)
;

\draw[gray,yshift=.5cm] (e1) 
arc [start angle=90, end angle=180, x radius={\df*1cm}, y radius={\df*1cm+.05cm}] 
arc [start angle=180, end angle=270, x radius={\df*1cm}, y radius={\df*1cm-.05cm}] 
;


%3
\def\dg{1.5}
\path[yshift=.25cm,->] (0.25,0) 
arc [start angle=-90, end angle=0, x radius={\dg*1cm}, y radius={\dg*1cm-.1cm}]
arc [start angle=0, end angle=90, x radius={\dg*1cm}, y radius={\dg*1cm+.1cm}]
coordinate (e1)
;

\draw[gray,yshift=.5cm] (e1) 
arc [start angle=90, end angle=180, x radius={\dg*1cm}, y radius={\dg*1cm+.05cm}] 
arc [start angle=180, end angle=270, x radius={\dg*1cm}, y radius={\dg*1cm-.05cm}] 
;

%reverse
\begin{scope}[yscale=-1,gray]

%1
\path[yshift=.75cm,->] (0.25,0) 
arc [start angle=-90, end angle=0, x radius={\ds*1cm}, y radius={\ds*1cm-.02cm}]
arc [start angle=0, end angle=90, x radius={\ds*1cm}, y radius={\ds*1cm+.05cm}]
coordinate (e1)
;

\draw[yshift=.75cm] (e1) 
arc [start angle=90, end angle=180, x radius={\ds*1cm}, y radius={\ds*1cm+.05cm}] 
arc [start angle=180, end angle=270, x radius={\ds*1cm}, y radius={\ds*1cm-.02cm}] 
;


%2
\def\df{.9}
\path[yshift=.5cm,->] (0.25,0) 
arc [start angle=-90, end angle=0, x radius={\df*1cm}, y radius={\df*1cm-.1cm}]
arc [start angle=0, end angle=90, x radius={\df*1cm}, y radius={\df*1cm+.1cm}]
coordinate (e1)
;

\draw[yshift=.5cm] (e1) 
arc [start angle=90, end angle=180, x radius={\df*1cm}, y radius={\df*1cm+.05cm}] 
arc [start angle=180, end angle=270, x radius={\df*1cm}, y radius={\df*1cm-.05cm}] 
;


%3
\def\dg{1.5}
\path[yshift=.25cm,->] (0.25,0) 
arc [start angle=-90, end angle=0, x radius={\dg*1cm}, y radius={\dg*1cm-.1cm}]
arc [start angle=0, end angle=90, x radius={\dg*1cm}, y radius={\dg*1cm+.1cm}]
coordinate (e1)
;

\draw[yshift=.5cm] (e1) 
arc [start angle=90, end angle=180, x radius={\dg*1cm}, y radius={\dg*1cm+.05cm}] 
arc [start angle=180, end angle=270, x radius={\dg*1cm}, y radius={\dg*1cm-.05cm}] 
;

\draw[decoration={markings,mark=at position {1/2+1.5/20} with {\arrow{<}}},postaction={decorate}] (.25,0) --++ (-2,0);

\end{scope}


%%%%%%%%%%%%%%%%%%%%%%%%%%%%%
%ring
\begin{scope}
    
\draw[line width=.4pt, decoration={width and color change,start color=brown, end color=brown}, decorate] (0.5,{.0pt})
arc [start angle={0}, end angle={180}, x radius={\xr*1cm}, y radius=1cm] 
% arc [start angle=270, end angle=180, x radius=.5cm, y radius=1cm] 
;

\draw[line width=.4pt, decoration={width and color change,start color=brown, end color=brown}, decorate] (0.5,{-.0pt})
% arc [start angle={90}, end angle={180}, x radius=.5cm, y radius=1cm] 
arc [start angle=360, end angle=180, x radius={\xr*1cm}, y radius=1cm] 
;

%arrows
\path[red,->,thin] (-0.2,0)
% arc [start angle={90}, end angle={180}, x radius=.5cm, y radius=1cm] 
arc [start angle=180, end angle=190, x radius={\xr*1cm}, y radius=1cm] coordinate (a) 
;

\draw[red,->,thin] (a)
% arc [start angle={90}, end angle={180}, x radius=.5cm, y radius=1cm] 
arc [start angle=190, end angle=260, x radius={\xr*1cm}, y radius=1cm] node[black,below right,shift={(-.16,.07)}] {$I$}
;

\end{scope}


%%%%%%%%%%%%%%%%%%%%%%%%%%%%%
%%%%%%%%%%%%%%%%%%%%%%%%%%%%%


%lines right
%1
\draw[gray,yshift=.75cm,->] (0.25,0) 
arc [start angle=-90, end angle=0, x radius={\ds*1cm}, y radius={\ds*1cm-.02cm}]
arc [start angle=0, end angle=90, x radius={\ds*1cm}, y radius={\ds*1cm+.05cm}]
coordinate (e1)
;


%2
\draw[gray,yshift=.5cm,->] (0.25,0) 
arc [start angle=-90, end angle=0, x radius={\df*1cm}, y radius={\df*1cm-.1cm}]
arc [start angle=0, end angle=90, x radius={\df*1cm}, y radius={\df*1cm+.1cm}]
coordinate (e1)
;


%3
\def\dg{1.5}
\draw[gray,yshift=.25cm,->] (0.25,0) 
arc [start angle=-90, end angle=0, x radius={\dg*1cm}, y radius={\dg*1cm-.1cm}]
arc [start angle=0, end angle=90, x radius={\dg*1cm}, y radius={\dg*1cm+.1cm}]
coordinate (e1)
;


%reverse
\begin{scope}[yscale=-1,gray]

%1
\draw[yshift=.75cm,->] (0.25,0) 
arc [start angle=-90, end angle=0, x radius={\ds*1cm}, y radius={\ds*1cm-.02cm}]
arc [start angle=0, end angle=90, x radius={\ds*1cm}, y radius={\ds*1cm+.05cm}]
coordinate (e1)
;


%2
\draw[yshift=.5cm,->] (0.25,0) 
arc [start angle=-90, end angle=0, x radius={\df*1cm}, y radius={\df*1cm-.1cm}]
arc [start angle=0, end angle=90, x radius={\df*1cm}, y radius={\df*1cm+.1cm}]
coordinate (e1)
;


%3
\def\dg{1.5}
\draw[yshift=.25cm,->] (0.25,0) 
arc [start angle=-90, end angle=0, x radius={\dg*1cm}, y radius={\dg*1cm-.1cm}]
arc [start angle=0, end angle=90, x radius={\dg*1cm}, y radius={\dg*1cm+.1cm}]
coordinate (e1)
;

\draw[decoration={markings,mark=at position {1/2+1.5/20} with {\arrow{>}}},postaction={decorate}] (.25,0) --++ (2,0);

\end{scope}

    
    \end{tikzpicture}
\caption{Линии магнитного поля кольцевого провода с током}
\label{fig:p149}
\end{figure}


% Через любую точку пространства вне кольца можно провести соответствующую линию его магнитного поля. С удалением от кольца густота линий поля убывает (убывает индукция магнитного поля).

Линии поля кольца как бы выходят из него со одной стороны и входят в него с другой стороны (внутри кольца линии замыкаются).

Для определения того, с какой стороны из кольца выходят линии его магнитного поля, можно воспользоваться следующим правилом.
%
\begin{law}
    \textbf{Правило правой руки.} Направление линий магнитного поля \emph{внутри} кольца укажет большой палец правой руки, если остальные согнутые пальцы (охватывающие кольцо) расположены так, что они указывают направление тока кольца.
\end{law}
%
Рисунок~\ref{fig:p150} иллюстрирует применение этого правила.  

    \begin{figure}[H]
    \centering
\includegraphics[]{el/37ring_pic.pdf}%
\caption{Правило правой руки}
\label{fig:p150}
\end{figure}


Соотнося рисунки~\ref{fig:p149} и~\ref{fig:p150}, можно видеть, что большой палец на рис.\,\ref{fig:p150} указывает направление линий магнитного поля внутри кольца.


















 
\end{document}